% Options for packages loaded elsewhere
% Options for packages loaded elsewhere
\PassOptionsToPackage{unicode}{hyperref}
\PassOptionsToPackage{hyphens}{url}
\PassOptionsToPackage{dvipsnames,svgnames,x11names}{xcolor}
%
\documentclass[
  letterpaper,
  DIV=11,
  numbers=noendperiod]{scrreprt}
\usepackage{xcolor}
\usepackage{amsmath,amssymb}
\setcounter{secnumdepth}{5}
\usepackage{iftex}
\ifPDFTeX
  \usepackage[T1]{fontenc}
  \usepackage[utf8]{inputenc}
  \usepackage{textcomp} % provide euro and other symbols
\else % if luatex or xetex
  \usepackage{unicode-math} % this also loads fontspec
  \defaultfontfeatures{Scale=MatchLowercase}
  \defaultfontfeatures[\rmfamily]{Ligatures=TeX,Scale=1}
\fi
\usepackage{lmodern}
\ifPDFTeX\else
  % xetex/luatex font selection
\fi
% Use upquote if available, for straight quotes in verbatim environments
\IfFileExists{upquote.sty}{\usepackage{upquote}}{}
\IfFileExists{microtype.sty}{% use microtype if available
  \usepackage[]{microtype}
  \UseMicrotypeSet[protrusion]{basicmath} % disable protrusion for tt fonts
}{}
\makeatletter
\@ifundefined{KOMAClassName}{% if non-KOMA class
  \IfFileExists{parskip.sty}{%
    \usepackage{parskip}
  }{% else
    \setlength{\parindent}{0pt}
    \setlength{\parskip}{6pt plus 2pt minus 1pt}}
}{% if KOMA class
  \KOMAoptions{parskip=half}}
\makeatother
% Make \paragraph and \subparagraph free-standing
\makeatletter
\ifx\paragraph\undefined\else
  \let\oldparagraph\paragraph
  \renewcommand{\paragraph}{
    \@ifstar
      \xxxParagraphStar
      \xxxParagraphNoStar
  }
  \newcommand{\xxxParagraphStar}[1]{\oldparagraph*{#1}\mbox{}}
  \newcommand{\xxxParagraphNoStar}[1]{\oldparagraph{#1}\mbox{}}
\fi
\ifx\subparagraph\undefined\else
  \let\oldsubparagraph\subparagraph
  \renewcommand{\subparagraph}{
    \@ifstar
      \xxxSubParagraphStar
      \xxxSubParagraphNoStar
  }
  \newcommand{\xxxSubParagraphStar}[1]{\oldsubparagraph*{#1}\mbox{}}
  \newcommand{\xxxSubParagraphNoStar}[1]{\oldsubparagraph{#1}\mbox{}}
\fi
\makeatother


\usepackage{longtable,booktabs,array}
\usepackage{calc} % for calculating minipage widths
% Correct order of tables after \paragraph or \subparagraph
\usepackage{etoolbox}
\makeatletter
\patchcmd\longtable{\par}{\if@noskipsec\mbox{}\fi\par}{}{}
\makeatother
% Allow footnotes in longtable head/foot
\IfFileExists{footnotehyper.sty}{\usepackage{footnotehyper}}{\usepackage{footnote}}
\makesavenoteenv{longtable}
\usepackage{graphicx}
\makeatletter
\newsavebox\pandoc@box
\newcommand*\pandocbounded[1]{% scales image to fit in text height/width
  \sbox\pandoc@box{#1}%
  \Gscale@div\@tempa{\textheight}{\dimexpr\ht\pandoc@box+\dp\pandoc@box\relax}%
  \Gscale@div\@tempb{\linewidth}{\wd\pandoc@box}%
  \ifdim\@tempb\p@<\@tempa\p@\let\@tempa\@tempb\fi% select the smaller of both
  \ifdim\@tempa\p@<\p@\scalebox{\@tempa}{\usebox\pandoc@box}%
  \else\usebox{\pandoc@box}%
  \fi%
}
% Set default figure placement to htbp
\def\fps@figure{htbp}
\makeatother


% definitions for citeproc citations
\NewDocumentCommand\citeproctext{}{}
\NewDocumentCommand\citeproc{mm}{%
  \begingroup\def\citeproctext{#2}\cite{#1}\endgroup}
\makeatletter
 % allow citations to break across lines
 \let\@cite@ofmt\@firstofone
 % avoid brackets around text for \cite:
 \def\@biblabel#1{}
 \def\@cite#1#2{{#1\if@tempswa , #2\fi}}
\makeatother
\newlength{\cslhangindent}
\setlength{\cslhangindent}{1.5em}
\newlength{\csllabelwidth}
\setlength{\csllabelwidth}{3em}
\newenvironment{CSLReferences}[2] % #1 hanging-indent, #2 entry-spacing
 {\begin{list}{}{%
  \setlength{\itemindent}{0pt}
  \setlength{\leftmargin}{0pt}
  \setlength{\parsep}{0pt}
  % turn on hanging indent if param 1 is 1
  \ifodd #1
   \setlength{\leftmargin}{\cslhangindent}
   \setlength{\itemindent}{-1\cslhangindent}
  \fi
  % set entry spacing
  \setlength{\itemsep}{#2\baselineskip}}}
 {\end{list}}
\usepackage{calc}
\newcommand{\CSLBlock}[1]{\hfill\break\parbox[t]{\linewidth}{\strut\ignorespaces#1\strut}}
\newcommand{\CSLLeftMargin}[1]{\parbox[t]{\csllabelwidth}{\strut#1\strut}}
\newcommand{\CSLRightInline}[1]{\parbox[t]{\linewidth - \csllabelwidth}{\strut#1\strut}}
\newcommand{\CSLIndent}[1]{\hspace{\cslhangindent}#1}



\setlength{\emergencystretch}{3em} % prevent overfull lines

\providecommand{\tightlist}{%
  \setlength{\itemsep}{0pt}\setlength{\parskip}{0pt}}



 


\KOMAoption{captions}{tableheading}
\makeatletter
\@ifpackageloaded{bookmark}{}{\usepackage{bookmark}}
\makeatother
\makeatletter
\@ifpackageloaded{caption}{}{\usepackage{caption}}
\AtBeginDocument{%
\ifdefined\contentsname
  \renewcommand*\contentsname{Table of contents}
\else
  \newcommand\contentsname{Table of contents}
\fi
\ifdefined\listfigurename
  \renewcommand*\listfigurename{List of Figures}
\else
  \newcommand\listfigurename{List of Figures}
\fi
\ifdefined\listtablename
  \renewcommand*\listtablename{List of Tables}
\else
  \newcommand\listtablename{List of Tables}
\fi
\ifdefined\figurename
  \renewcommand*\figurename{Figure}
\else
  \newcommand\figurename{Figure}
\fi
\ifdefined\tablename
  \renewcommand*\tablename{Table}
\else
  \newcommand\tablename{Table}
\fi
}
\@ifpackageloaded{float}{}{\usepackage{float}}
\floatstyle{ruled}
\@ifundefined{c@chapter}{\newfloat{codelisting}{h}{lop}}{\newfloat{codelisting}{h}{lop}[chapter]}
\floatname{codelisting}{Listing}
\newcommand*\listoflistings{\listof{codelisting}{List of Listings}}
\makeatother
\makeatletter
\makeatother
\makeatletter
\@ifpackageloaded{caption}{}{\usepackage{caption}}
\@ifpackageloaded{subcaption}{}{\usepackage{subcaption}}
\makeatother
\usepackage{bookmark}
\IfFileExists{xurl.sty}{\usepackage{xurl}}{} % add URL line breaks if available
\urlstyle{same}
\hypersetup{
  pdftitle={.},
  pdfauthor={Norah Jones},
  colorlinks=true,
  linkcolor={blue},
  filecolor={Maroon},
  citecolor={Blue},
  urlcolor={Blue},
  pdfcreator={LaTeX via pandoc}}


\title{.}
\author{Norah Jones}
\date{2026-03-13}
\begin{document}
\maketitle

\renewcommand*\contentsname{Table of contents}
{
\hypersetup{linkcolor=}
\setcounter{tocdepth}{2}
\tableofcontents
}

\bookmarksetup{startatroot}

\chapter*{Preface}\label{preface}
\addcontentsline{toc}{chapter}{Preface}

\markboth{Preface}{Preface}

This is a Quarto book.

To learn more about Quarto books visit
\url{https://quarto.org/docs/books}.

\bookmarksetup{startatroot}

\chapter{Introduction}\label{introduction}

This is a book created from markdown and executable code.

See Knuth (1984) for additional discussion of literate programming.

\bookmarksetup{startatroot}

\chapter{项目名称}\label{ux9879ux76eeux540dux79f0}

银龄守望------基于人工智能的农村空巢老人安全监护与社区联动响应系统

\section{一、项目背景}\label{ux4e00ux9879ux76eeux80ccux666f}

农村人口空心化加剧，大量空巢老人独居，突发健康风险（如跌倒、失能）难以及时发现。传统人工巡访频次低、响应滞后，而老人又普遍缺乏主动求助能力。\textbf{人工智能作为新一代生产力工具，为破解这一``看不见的风险''提供了技术可能}。

\section{二、项目目标}\label{ux4e8cux9879ux76eeux76eeux6807}

\textbf{依托轻量化、高隐私的人工智能感知能力}，构建一套面向农村场景的智能监护与社区响应机制，实现对空巢老人异常状态的自动识别、分级预警与快速处置，推动基层养老服务从``被动响应''迈向``主动守护''。

\section{三、核心内容}\label{ux4e09ux6838ux5fc3ux5185ux5bb9}

\begin{itemize}
\tightlist
\item
  在老人家中部署\textbf{基于AI的行为识别设备}，无需穿戴、无需操作，\textbf{自动识别跌倒、长时间静止、夜间异常活动等高风险行为}；\\
\item
  \textbf{AI系统实时分析并触发分级预警}，通过微信、短信或电话将信息精准推送至村委、网格员或亲属；\\
\item
  建立``AI监测---社区响应---家庭协同''的闭环机制，\textbf{将AI能力深度嵌入村级治理末梢}。
\end{itemize}

\section{四、科学性、先进性与独立性}\label{ux56dbux79d1ux5b66ux6027ux5148ux8fdbux6027ux4e0eux72ecux7acbux6027}

\subsection{科学性}\label{ux79d1ux5b66ux6027}

项目基于老年行为医学与社会照护理论，\textbf{AI识别规则由医学专家与基层工作者共同校准}，确保算法判断符合真实生活场景（如区分``躺卧休息''与``跌倒昏迷''），避免误报漏报。

\subsection{先进性}\label{ux5148ux8fdbux6027}

\textbf{本项目的核心创新在于将人工智能从城市实验室下沉到乡村家庭}，首次实现：\\
- \textbf{无感化AI监护}：老人无需学习、操作，技术``隐形''融入生活；\\
- \textbf{边缘化AI部署}：模型运行于本地设备，保障隐私、适应弱网环境；\\
-
\textbf{治理化AI应用}：AI不仅是工具，更是连接``家庭---社区---政府''的智能纽带。\\
这代表了AI在公共服务领域从``炫技''走向``务实''的重要转向。

\subsection{独立性}\label{ux72ecux7acbux6027}

系统\textbf{不依赖云端算力或专业运维团队}，AI模型轻量化、本地化运行，村级即可完成安装与基础管理；预警逻辑可按村调整，\textbf{真正实现``AI为村所用、由村主导''}，具备高度自主性和推广韧性。

\section{五、现实利用价值与现实意义}\label{ux4e94ux73b0ux5b9eux5229ux7528ux4ef7ux503cux4e0eux73b0ux5b9eux610fux4e49}

\subsection{现实利用价值}\label{ux73b0ux5b9eux5229ux7528ux4ef7ux503c}

\begin{itemize}
\tightlist
\item
  \textbf{对政府}：以较低成本引入AI能力，显著提升基层养老兜底效率，\textbf{用技术弥补人力缺口}；\\
\item
  \textbf{对村委/网格员}：AI承担``7×24小时盯守''任务，\textbf{释放人力聚焦情感关怀与复杂处置}；\\
\item
  \textbf{对家庭}：子女获得``数字安心''，\textbf{AI成为远方亲情的延伸}；\\
\item
  \textbf{对老人}：在尊严与隐私受保护的前提下，获得全天候安全保障。
\end{itemize}

\subsection{现实意义}\label{ux73b0ux5b9eux610fux4e49}

在老龄化与数字化双重浪潮下，\textbf{AI不应只服务于城市精英，更应惠及最脆弱的农村老人}。本项目探索了一条``\textbf{普惠型AI+基层治理}''的中国路径，证明人工智能可以成为缩小城乡服务差距、实现社会公平的重要力量，具有鲜明的时代价值与政策示范意义。

\section{六、赛道契合}\label{ux516dux8d5bux9053ux5951ux5408}

\begin{itemize}
\tightlist
\item
  \textbf{主赛道}：现代服务与社会治理（AI赋能社区服务与基层治理）\\
\item
  \textbf{辅助赛道}：新一代信息技术（AI为核心支撑）、生物医药与健康科技（AI驱动健康监护）
\end{itemize}

\section{七、团队与资源}\label{ux4e03ux56e2ux961fux4e0eux8d44ux6e90}

\begin{itemize}
\tightlist
\item
  \textbf{核心团队}：兼具AI应用开发、公共政策与社会创新经验，确保技术``可用、可信、可持续''；\\
\item
  \textbf{政府合作}：地方民政部门支持将\textbf{AI监护纳入村级智慧治理标配}；\\
\item
  \textbf{学术支撑}：高校团队协助验证AI行为识别的准确性与伦理边界；\\
\item
  \textbf{社区协同}：村干部参与AI预警规则本地化调试，\textbf{让算法``听得懂方言、认得清乡情''}；\\
\item
  \textbf{产业支持}：国产AI芯片厂商提供适配农村场景的边缘计算硬件。
\end{itemize}

\section{八、实施路径}\label{ux516bux5b9eux65bdux8defux5f84}

\begin{enumerate}
\def\labelenumi{\arabic{enumi}.}
\tightlist
\item
  试点阶段：在行政村部署AI监护终端，验证\textbf{AI识别准确率与社区响应时效}；\\
\item
  优化阶段：根据反馈调整AI模型与预警阈值，\textbf{提升``技术---治理''协同效率}；\\
\item
  推广阶段：推动将``AI助老''纳入县域民生工程，\textbf{实现从``样板间''到``普及房''跨越}。
\end{enumerate}

\begin{center}\rule{0.5\linewidth}{0.5pt}\end{center}

本项目终围绕一个核心逻辑：\textbf{AI不是炫技，而是解决真问题的关键使能器}。它在农村养老场景中实现了``看得见风险、连得上响应、守得住底线''的独特价值，且完全服务于社会治理目标。

\bookmarksetup{startatroot}

\chapter*{References}\label{references}
\addcontentsline{toc}{chapter}{References}

\markboth{References}{References}

\phantomsection\label{refs}
\begin{CSLReferences}{1}{0}
\bibitem[\citeproctext]{ref-knuth84}
Knuth, Donald E. 1984. {``Literate Programming.''} \emph{Comput. J.} 27
(2): 97--111. \url{https://doi.org/10.1093/comjnl/27.2.97}.

\end{CSLReferences}




\end{document}
